\section{Conclusion}
\label{sec:conclusion}

This paper presented a PRISMA 2020 systematic review and bibliometric analysis of Generative Artificial Intelligence and Large Language Models in digital forensics, synthesizing $N=44$ peer-reviewed studies (22 journal articles and 22 archival proceedings) retrieved via 28 Crossref REST API queries from 2021 through November 10, 2026 and cross-validated via OpenAlex (44/44 DOI match; 97.7\% affiliation and 90.9\% abstract/full-text coverage).

Quantitative bibliometrics (RQ1) confirm explosive recent expansion, with 84.1\% ($n=37$) of the corpus published in 2025--2026 year-to-date and \textit{Forensic Science International: Digital Investigation} leading dissemination at 29.5\% ($n=13$). Decoupled orthogonal coding resolves prior taxonomic ambiguities by distinguishing LLMs applied \textit{to} DFIR ($n=39$, 88.6\%) from the emerging discipline of digital forensics \textit{of} LLM and agentic platforms ($n=5$, 11.4\%). Operational synthesis (RQ2) confirms concrete triage and classification gains across report generation, knowledge-graph verification, behavioral malware detection, volatile memory forensics, and multi-agent incident response, while exposing a pervasive reproducibility deficit: zero of six audited open-access empirical full texts disclosed complete execution manifests.

Legal and cryptographic analysis (RQ3) together with adversarial evaluation (RQ4) identify formidable barriers to judicial adoption. Among 34 empirical LLM-for-DFIR tool studies, 67.6\% ($n=23$) rely on commercial cloud APIs or omit local execution guarantees (versus 23.5\%, $n=8$, dedicated local/offline models and 8.8\%, $n=3$, cross-model open-weight benchmarks), conflicting directly with GDPR, the EU Law Enforcement Directive (2016/680) and Portuguese Law n.{\textordmasculine}~59/2019, judicial secrecy under CPP Article~86, and high-risk AI mandates under Annex~III of the EU AI Act (Regulation 2024/1689). Simultaneously, GPU floating-point non-associativity disrupts exact replay, uncritical acceptance of unsupported premises exposes prompt-dependent error rates under FRE 702(d), and evidence-borne prompt injection turns untrusted artifacts into active control-flow attacks.

Language models cannot replace deterministic forensic tooling or sworn human expert testimony. Their legitimate role in digital forensics is strictly bounded to air-gapped cognitive coprocessors anchored by ISO/IEC 27037/27041/27042 controls, deterministic-first parsing, cryptographic execution manifests, and personal, non-delegable examiner accountability.
